\documentclass[../../../include/open-logic-section]{subfiles}

\begin{document}

\olfileid{sth}{cardinals}{hp}
\olsection{ضمیمہ: ہیوم دا اصول}

\olref[cp]{sec} وچ اسیں کانتور دا اصول دَسیا سی۔ اوہ ایہ سی:
\begin{align*}
	\card{A} = \card{B} & \text{ اُدوں تے صرف اُدوں جدوں } A \approx B.
\intertext{ایہ اوہدی طرح دا ای اے جیہنوں ہُن \emph{ہیوم دا اصول} آکھدے نیں۔ اوہ ایہ آکھدا اے:}
	\fregenum{x} {F(x)} = \fregenum{x}{G(x)} & \text{ اُدوں تے صرف اُدوں جدوں } F \sim G
\end{align*}
جتھے «$F \sim G$» دا مختصر مطلب ایہ اے کہ $F$ شےواں این اونیاں ای نیں جِنّیاں $G$ شےواں، یعنی $F$ شےواں نوں $G$ شےواں نال اک دوطرفی نقشہ بندی وچ رکھ سکدے آں، یعنی:
\begin{align*}
	\exists R(&\forall v\forall y(Rvy \lif (Fv \land Gy)) \land {}\\
		&\forall v(Fv \lif \lexists![y][Rvy]) \land {}\\
		&\forall y(Gy \lif \lexists![v][Rvy]))
\end{align*}
پر ہیوم دے اصول تے کانتور دے اصول وچ نوع دا فرق اے۔ کانتور دے اصول دے بیان وچ متغیر «$A$» تے «$B$» پہلے درجے دیاں ٹرماں نیں جیہڑیاں \emph{سیٹاں} لئی کھڑیاں نیں۔ ہیوم دے اصول دے بیان وچ «$F$»، «$G$» تے «$R$» پہلے درجے دیاں ٹرماں \emph{نہیں}۔ اوہ \emph{محمول دی تھاں} نیں۔ (شاید اوہ \emph{خاصیتاں} لئی کھڑے نیں۔) سو ہیوم دے اصول نوں عام بولی وچ ایویں آکھ سکدے آں۔ $F$ شےواں دی گنتی $G$ شےواں دی گنتی اے اُدوں تے صرف اُدوں جدوں $F$ شےواں تے~$G$ شےواں وچ دوطرفی نقشہ بندی اے۔ ایس نوں \emph{ہیوم دا اصول} آکھدے نیں، کیوں جے ہیوم نے اک واری ایہ لکھیا:
\begin{quote}
  جدوں دو عدد ایس طرح جوڑے جان کہ اک دے ہر اک جُز دے سامنے دوجے دا اک جُز ہووے، تاں اسیں اوہناں نوں برابر آکھدے آں۔
  \citep[Pt.III Bk.1 \S1]{Hume1740}
\end{quote}
تے اج دی ریاضیاتی منطق وچ ہیوم دے اصول نوں اہم تھاں \citet[\S63]{Frege1884} نے دِوائی۔ اوہنے ہیوم دی ایہ عبارت نقل کیتی تے فیر (عملی طور اُتے) اوہدا خاکہ دتا جیہنوں اسیں ہیوم دا اصول آکھیا اے۔

ہیوم دے اصول تے پنجویں بنیادی قانون دی ساخت وچ مشابہت ول دھیان دیو۔ اسیں اوہ قانون \olref[story][blv]{sec} وچ ایویں لکھیا سی:
\[
	\fregeext{x}{F(x)} = \fregeext{x}{G(x)} \text{ اُدوں تے صرف اُدوں جدوں } \lforall[x][(F(x) \liff G(x))].
\]
ایتھے کجھ وضاحت تے مقابلہ کم آ سکدے نیں۔

ہیوم دے اصول یا پنجویں بنیادی قانون ورگے اصول نوں لین دے دو طریقے نیں، \emph{اسنادی طور اُتے} یا \emph{غیر اسنادی طور اُتے} (\olref[story][predicative]{sec} چیتے کرو)۔ پنجویں بنیادی قانون دی غیر اسنادی پڑھت وچ، ہر~$F$ لئی، شے $\fregeext{x}{F(x)}$ اوہناں مسوّراں دے دائرے وچ آؤندی اے جیہڑے اسیں پنجواں بنیادی قانون آپ بیان کرن لئی ورتے۔ ایویں ای ہیوم دے اصول دی غیر اسنادی پڑھت وچ، ہر~$F$ لئی، شے $\fregenum{x}{F(x)}$ اوہناں مسوّراں دے دائرے وچ آؤندی اے جیہڑے اسیں ہیوم دا اصول بیان کرن لئی ورتے۔ دوجے پاسے، \emph{اسنادی} سمجھ مطابق، شےواں $\fregeext{x}{F(x)}$ تے~$\fregenum{x}{F(x)}$ کسے \emph{وکھرے} دائرے دیاں شےواں ہون گیاں۔

ہُن جے اسیں پنجواں بنیادی قانون غیر اسنادی طور اُتے پڑھیے، تاں سادہ خاصیت توں سیٹ بناؤن دے اصول راہیں ٹکراؤ نکلدا اے (تفصیل لئی \olref[story][blv]{sec} ویکھو)۔ سادہ خاصیت توں سیٹ بناؤن دے اصول وانگ، اینوں \emph{اسنادی طور اُتے} پڑھ کے ٹکراؤ توں بچایا جا سکدا اے۔ پر شاید ایہ اوہ سارا کم نہیں کرے گا جیہڑا اسیں چاہندے ساں۔

پر ہیوم دا اصول غیر اسنادی طور اُتے بے ٹکراؤ پڑھیا جا \emph{سکدا} اے۔ تے ایس پڑھت وچ ایہ خاصا طاقتور اے۔

مثال لئی محمول «$x \neq x$» ویکھو، جیہدے اُتے صاف اے کہ کوئی شے پوری نہیں اُتر دی۔ ہُن ہیوم دے اصول توں شے $\# x( x\neq
x)$ ملدی اے۔ اسیں اینوں عدد~$0$ منّ سکدے آں۔ ہُن \emph{غیر اسنادی} سمجھ مطابق، تے \emph{صرف} غیر اسنادی سمجھ مطابق، ایہ شے $0$ ساڈے اصل مسوّراں دے دائرے وچ آؤندی اے۔ سو محمول «$x = 0$» نال ہیوم دا اصول لانا معنی رکھدا اے، تے ایویں سانوں شے $\#x (x =
0)$ ملدی اے۔ اسیں اینوں عدد~$1$ منّ سکدے آں۔ نالے ہیوم دے اصول توں لازم آؤندا اے کہ $0 \neq 1$، کیوں جے اپنے آپ نال اکّو شے نہ ہون والیاں شےواں توں $0$ نال اکّو شے ہون والیاں شےواں ول دوطرفی نقشہ بندی نہیں ہو سکدی (پہلی قسم وچ کوئی شے نہیں، دوجی وچ اک اے)۔ ہُن فیر غیر اسنادی طور اُتے کم کردیاں، $1$~ساڈے اصل مسوّراں دے دائرے وچ آؤندا اے۔ سو محمول «$(x = 0 \lor x = 1)$» نال ہیوم دا اصول لانا معنی رکھدا اے، تے ایویں سانوں شے $\#x(x = 0 \lor
x = 1)$ ملدی اے۔ اسیں اینوں عدد~$2$ منّ سکدے آں، تے وکھا سکدے آں کہ $0\neq 2$ تے $1 \neq 2$، تے ایس طرح اگے۔

سو مختصر گل ایہ اے کہ غیر اسنادی طور اُتے لئیے تاں ہیوم دے اصول توں لازم آؤندا اے کہ شےواں \emph{لامتناہی گنتی وچ} نیں۔ ایس گل نے \emph{نویں فریگے والے منطقیت پسنداں} نوں حساب دی بنیاد وجوں ہیوم دا اصول لین ول مائل کیتا اے۔

پر فریگے نے \emph{آپ} ہیوم دا اصول اپنے حساب دی بنیاد نہیں منّیا۔ اوہنے ہیوم دا اصول اک صاف تعریف توں ثابت کیتا۔ $\fregenum{x}{F(x)}$ دی تعریف تصور $F \sim \Phi$ دے دائرے وجوں دتی اے۔ اج دی لکھت وچ اسیں ایہنوں $\fregenum{x}{F(x)} = \Setabs{G}{F \sim G}$ وجوں لکھن دی کوشش کر سکدے آں، پر ایہ سانوں سادہ خاصیت توں سیٹ بناؤن دے اصول دیاں اوکڑاں ول موڑ لے جائے گا۔

\end{document}
